\begin{table}[htbp]
  \centering
  \caption{Mean time for one simulation step: the scene total, median over repeats, divided by the number of frames. A step runs both query types, so this is the cost of the vertex-face and edge-edge work of that frame together. \emph{BP full} is the whole broad phase, the acceleration structure plus the queries over it; \emph{BP prep} is the structure alone and \emph{BP queries} the queries alone. The narrow phase is named by the answer it is asked for: \emph{NP EToI} returns one earliest time of impact for the step, so every query prunes against the running minimum, and \emph{NP per-pair} returns one per candidate with no shared bound. BP full here includes building the swept boxes.}
  \label{tab:per-frame}
  \fittable{%
  \begin{tabular}{lrlrrr}
    \toprule
    scene & frames & mode & BP full (ms) & NP EToI (ms) & total ms \\
    \midrule
    armadillo-rollers & 396 & GPU & 3.03 & 3.91 & 6.94 \\
    armadillo-rollers & 396 & CPU & 3.48 & 1.74 & 5.22 \\
    cloth-ball & 43 & GPU & 15.10 & 4.18 & 19.28 \\
    cloth-ball & 43 & CPU & 13.23 & 14.39 & 27.62 \\
    cloth-funnel & 372 & GPU & 4.20 & 2.10 & 6.30 \\
    cloth-funnel & 372 & CPU & 2.36 & 1.18 & 3.55 \\
    n-body & 74 & GPU & 52.61 & 6.43 & 59.04 \\
    n-body & 74 & CPU & 42.60 & 76.48 & 119.08 \\
    puffer-ball & 120 & GPU & 175.35 & 11.10 & 186.45 \\
    puffer-ball & 120 & CPU & 206.72 & 175.38 & 382.10 \\
    rod-twist & 2,556 & GPU & 3.45 & 6.48 & 9.93 \\
    rod-twist & 2,556 & CPU & 6.32 & 12.41 & 18.72 \\
    \bottomrule
  \end{tabular}}
  \par\smallskip\footnotesize A mean rather than a median over steps: the scene total is what a run costs, and the mean is the only average that divides back into it.
  \par\smallskip\footnotesize Source: \texttt{benchmark/assessment/broadphase-cell2dmin.csv}
\end{table}
